\section{Introduction}
\label{sec:intro}

Multi-task learning (MTL) is a central pillar of modern artificial intelligence, enabling a single model to generalize across a diverse spectrum of domains and applications. While traditional multi-task training paradigms require simultaneous joint optimization over large-scale aggregated datasets, they suffer from extreme computational costs, data privacy concerns, and severe negative transfer (or task interference). To circumvent these bottlenecks, weight-space post-hoc model merging has emerged as a groundbreaking, training-free alternative. This paradigm allows separate expert models, independently fine-tuned from a shared pretrained initialization, to be directly composed into a single multi-task network without accessing the original training datasets.

Early model merging frameworks, such as \textbf{Task Arithmetic} \cite{task_vectors} and \textbf{TIES-Merging} \cite{ties_merging}, employ uniform, hand-tuned scaling coefficients to blend model weights. While simple and computationally lightweight, these static schemes ignore the complex, layer-specific correlations and conflicts between tasks, often resulting in sub-optimal joint performance. To resolve this, \textbf{Adaptive Model Merging (AdaMerging)} \cite{adamerging} introduced test-time adaptation (TTA), which dynamically learns layer-wise or task-wise merging coefficients by minimizing unsupervised entropy over a small calibration stream of unlabeled test queries.

\begin{figure}[t]
    \centering
    \includegraphics[width=0.9\linewidth]{comparison_plot.png}
    \caption{Model merging performance across different post-training quantization (PTQ) schemas. While unregularized test-time adaptation (AdaMerging) achieves peak performance in FP32, it experiences catastrophic representation collapse under INT8 and INT4 quantization. Our proposed Hessian-Regularized Coefficient Optimization (HessMerge) flattens the local loss landscape, delivering robust performance across all deployment quantization schemas.}
    \label{fig:teaser}
\end{figure}

Despite the empirical success of TTA in continuous floating-point (FP32) spaces, we identify a critical, hitherto overlooked vulnerability in these adaptive paradigms: \textbf{Quantization-Operator Overfitting}. In real-world edge deployment, merged models are routinely quantized to low-precision formats (e.g., INT8 or INT4 symmetric or asymmetric schemas) to meet stringent hardware memory and latency constraints. We find that unregularized TTA optimizes the merging coefficients into extremely sharp local minima. While these sharp minima yield high performance in FP32, they are highly sensitive to even minor weight perturbations. When subjected to post-training quantization (PTQ), the resulting rounding noise triggers a catastrophic collapse in multi-task accuracy. This represents a fundamental hurdle for deploying test-time merged models on edge devices.

To neutralize this vulnerability, we propose \textbf{Hessian-Regularized Coefficient Optimization (HessMerge)}, a mathematically rigorous framework that guarantees post-training quantization robustness for adaptive model merging. HessMerge formulates the quantization process as a bounded perturbation in parameter space. Through a second-order Taylor expansion, we show that the quantization-induced loss gap is directly bounded by the curvature of the loss landscape with respect to the merging coefficients. HessMerge explicitly incorporates a regularizer penalizing the \textbf{trace of the Hessian} ($\text{Tr}(\mathcal{H}_{\Theta})$) with respect to the continuous blending coefficients during test-time adaptation. By minimizing this trace, HessMerge flattens the local loss landscape, mathematically ensuring that weight-space perturbations from any downstream quantization operator result in a minimal drop in performance.

Our main contributions are summarized as follows:
\begin{itemize}
    \item We identify and analyze the phenomenon of \textbf{Quantization-Operator Overfitting} in test-time adaptive model merging, showing that unregularized coefficient optimization converges to sharp minima that collapse under post-training quantization.
    \item We propose \textbf{HessMerge}, a novel, mathematically rigorous framework that incorporates exact Hessian trace regularization into coefficient optimization to flatten the local loss landscape.
    \item We conduct a comprehensive multi-schema evaluation (including FP32, INT8 symmetric/asymmetric tensor-wise and channel-wise, and INT4 symmetric channel-wise quantization) across four diverse visual domain datasets (MNIST, FashionMNIST, CIFAR-10, SVHN).
    \item We demonstrate that HessMerge consistently outperforms existing state-of-the-art merging baselines, preserving high multi-task performance even under aggressive INT4 quantization and establishing a new frontier for robust model merging.
\end{itemize}
